\section{LoRA 用低秩把微调成本压下来}
\label{sec:14-Deep-Learning/09-Scaling-and-Adaptation/03}

一个七十亿参数的基座模型要服务四十个客户，每个客户的语料、术语和格式偏好都不一样。按常规做法，每个客户各做一次全参数微调，于是磁盘上躺着四十份七十亿参数的副本，切换客户意味着换掉整个模型的权重。存储是一笔钱，显存里同时装不下几份是另一笔，把一次请求路由到正确的副本又是第三笔。

真正被浪费掉的是一个显而易见的事实：这四十份权重彼此之间差别很小。它们都是从同一个 \(W_0\) 出发，各自挪了一小步。既然差异才是有信息的部分，为什么要存四十份完整的模型？

\subsection{只学增量，并且限制它的秩}

把微调后的权重写成基座加增量：

\[
  W = W_0 + \Delta W .
\]

\(W_0\) 冻结，不参与梯度更新。要省钱，就得让 \(\Delta W\) 比 \(W_0\) 便宜得多，而对一个 \(d \times k\) 的矩阵，最直接的省法是限制它的秩。LoRA 的做法是把增量写成两个瘦矩阵的乘积：

\[
  W = W_0 + BA,
  \qquad B \in \R^{d \times r},\quad A \in \R^{r \times k},\quad r \ll \min(d, k),
\]

训练时只更新 \(A\) 和 \(B\)。参数量从 \(dk\) 变成 \(r(d + k)\)。

这个比值下降得比直觉快。取 \(d = k = 4096\)、\(r = 8\)：\(dk \approx 1.7 \times 10^{7}\)，而 \(r(d+k) = 8 \times 8192 \approx 6.6 \times 10^{4}\)，比值约 \(1/256\)，也就是不到 \(0.4\%\)。四十个客户的全部差异加起来，还不到一份基座模型的两成。

\begin{center}
\begin{tikzpicture}[scale=1.0]

% W0
\draw[fill=black!12] (0,0) rectangle (3.4,3.4);
\node at (1.7,1.9) {\(W_0\)};
\node at (1.7,1.45) {\footnotesize 冻结};
\draw[<->,>=stealth] (-0.45,0) -- (-0.45,3.4);
\node[anchor=east] at (-0.5,1.7) {\footnotesize \(d\)};
\draw[<->,>=stealth] (0,-0.4) -- (3.4,-0.4);
\node[anchor=north] at (1.7,-0.45) {\footnotesize \(k\)};

\node at (3.85,1.7) {\large \(+\)};

% B
\draw[fill=black!45] (4.68,0) rectangle (5.03,3.4);
\node[anchor=south] at (4.855,3.45) {\footnotesize \(B\)};
\draw[<->,>=stealth] (4.68,-0.4) -- (5.03,-0.4);
\node[anchor=north] at (4.855,-0.45) {\footnotesize \(r\)};

\node at (5.5,1.7) {\large \(\times\)};

% A
\draw[fill=black!45] (6.0,1.525) rectangle (9.4,1.875);
\node[anchor=south] at (7.7,1.93) {\footnotesize \(A\)};
\draw[<->,>=stealth] (6.0,1.15) -- (9.4,1.15);
\node[anchor=north] at (7.7,1.10) {\footnotesize \(k\)};
\draw[<->,>=stealth] (9.7,1.525) -- (9.7,1.875);
\node[anchor=west] at (9.75,1.7) {\footnotesize \(r\)};

% 数字对照
\node[anchor=north,align=center,font=\footnotesize] at (1.7,-0.95)
  {\(dk = 4096 \times 4096\)\\ \(\approx 1.7 \times 10^{7}\)};
\node[anchor=north,align=center,font=\footnotesize] at (7.2,-0.95)
  {\(r(d+k) = 8 \times 8192\)\\ \(\approx 6.6 \times 10^{4}\)（约 \(0.4\%\)）};

\draw[dashed] (4.35,-2.0) -- (4.35,3.95);

\end{tikzpicture}
\end{center}

\emph{左边的方块是冻结的 \(W_0\)，右边一条竖细条乘一条横细条就是全部可训练参数。面积之比就是参数量之比，而图上那两条细条为了看得见已经画粗了——按真实比例它们细到画不出来，数值上只占 \(0.4\%\)。要记住的是这张图的形状——省钱来自 \(r\) 同时压缩了两个方向，而不是来自任何近似技巧。}

\subsection{它和低秩逼近是两件事}

看到``秩 \(r\)''和``矩阵分解''，很容易联想到 SVD，这个联想一半对一半错，值得拆开。

对的那一半是背景知识。任意矩阵都有奇异值分解（\hyperref[sec:03-Linear-Algebra/07-SVD-and-Data-Applications/01]{``任意矩阵的奇异值分解''}），而 Eckart--Young 定理说：在 Frobenius 范数或谱范数下，矩阵 \(M\) 的最优秩-\(r\) 逼近就是把 SVD 里最大的 \(r\) 个奇异值留下、其余截断得到的那个矩阵，逼近误差由被丢掉的奇异值决定（\hyperref[sec:03-Linear-Algebra/07-SVD-and-Data-Applications/03]{``低秩近似：用少数方向保留主要结构''}）。这告诉我们``秩 \(r\) 能装下多少信息''这个问题有明确答案，也告诉我们奇异值谱的衰减快慢决定了低秩表示划不划算。

错的那一半是把 LoRA 当成在做这件事。Eckart--Young 处理的是一个\emph{已知}矩阵：\(M\) 摆在那里，求离它最近的低秩矩阵，问题是逼近。LoRA 面对的是一个\emph{未知}矩阵：那个理想的 \(\Delta W\) 从来没有被算出来过，因为算它就等于做一次全参数微调，正是要避免的事。LoRA 做的是在秩不超过 \(r\) 的矩阵集合上直接跑梯度下降，最小化下游损失，问题是优化。

两者的差别不只是措辞。逼近问题有闭式解，最优性由定理保证；优化问题的可行集 \(\{\,M : \rank(M) \le r\,\}\) 不是凸集，而且用 \(BA\) 这个乘积形式参数化之后，即使原本关于 \(\Delta W\) 是凸的目标也会变成非凸的。训练结束时拿到的是这个非凸问题的一个驻点，不是``最优的秩-\(r\) 增量''。把 Eckart--Young 的最优性挪过来给 LoRA 背书，是一次不成立的搬运。

顺带说一句实践中的联系：\(\Delta W\) 的\emph{数值}秩——奇异值中显著大于噪声水平的那些的个数——才是判断低秩表示是否够用的量，而数值秩与代数秩往往差很远（\hyperref[sec:08-Numerical-Analysis/05-Numerical-Linear-Algebra/06]{``SVD、低秩与数值秩''}）。如果事后把某次全参数微调的 \(\Delta W\) 拿去做 SVD，看到的是奇异值缓慢衰减而非断崖，那就说明这个任务上的低秩假设不牢。

\subsection{那个假设是经验的}

LoRA 能成立，靠的是一个前提：微调所需要的权重改变量本身接近低秩。这是经验观察，不是任何定理的推论，而且它成立的程度随任务变化。

有一种直觉支持它：微调通常不是教模型新知识，而是调整它已有能力的调用方式——换一种语气、遵守一种格式、在几个领域术语上重新加权。这类调整落在一个远比全空间小的子空间里，是说得通的。但直觉在这里只能到这一步，它不能预测某个具体任务需要多大的 \(r\)。

可以操作的判断是这样的。\(r\) 太小时症状是欠拟合：训练损失下不去，而且加数据、加轮数都不管用——真正的增量放不进那个 \(r\) 维子空间，问题不在优化而在容量。\(r\) 增大时收益上升，但过了某个点后迅速递减，继续加 \(r\) 只是在给零奇异值方向配参数。所以合理的做法是从小的 \(r\) 起步，翻倍着往上试，观察训练损失的下降是否随之改善；一旦改善停止，就停在那里。风格与格式类的任务通常几个到十几的 \(r\) 就够，需要真正引入新领域知识的任务会明显更贪。

还有一处容易被忽略：不同层、不同投影矩阵对 \(r\) 的需求并不一致，把同一个 \(r\) 铺满全网络是简化而非最优。至于该给哪些矩阵加 LoRA，目前也没有原理性的答案，是靠试的。

\subsection{两个让它好用的细节}

\(B\) 通常初始化为零，\(A\) 用常规的随机初始化。于是训练第一步时 \(BA = 0\)，\(W = W_0\)，模型行为与基座完全一致。这不是可有可无的讲究：如果两个因子都随机初始化，训练一开始就给每个权重矩阵注入了一个随机扰动，基座模型的能力会先被打散再重新长回来，而这次打散是没有必要的。零初始化让微调成为从基座出发的连续形变，起点处损失就是基座的损失，之后单调地往下走。

另一个细节是推理时的合并。训练完成后可以直接算出 \(W = W_0 + BA\) 并存下来，之后的前向传播和一个普通模型没有区别，没有额外的矩阵乘法，没有额外的延迟。这一点把 LoRA 和许多别的高效微调方法区分开——那些往网络里插入额外模块的做法，插进去的东西在推理时也得跑，省下的是训练成本，付出的是永久的推理开销。LoRA 可以两头都省，代价只是合并之后就不能再免费地切换任务了：要保留``一个基座配多套适配器''的灵活性，就得让 \(BA\) 保持分离，此时那点额外的乘法又回来了。这是一个明确的取舍，按部署形态选。

实现里常见的 \(BA\) 前面乘一个 \(\alpha / r\) 形式的缩放常数，作用是让改变 \(r\) 时不必重新调学习率，属于工程约定，不影响上面任何一条判断。

\medskip

LoRA 改变的是模型的行为，成本压在训练侧，而模型的尺寸一点没变——合并之后它还是那个七十亿参数的模型，每生成一个 token 仍要把全部权重读一遍。想让它在部署时更便宜，就得动权重本身的表示精度，或者干脆换一个更小的模型来顶替它。这两条路各有各的代价，下一节把它们放在一起算账。
